\documentclass[11pt]{article}
\usepackage[a4paper,left=30mm,right=30mm,top=25mm,bottom=25mm,marginparwidth=16mm,marginparsep=4mm]{geometry}
\usepackage{fontspec,xcolor,pdflscape}
\usepackage[unicode,hidelinks]{hyperref}
\usepackage{polyglossia}
\setmainlanguage{english}
\setotherlanguages{arabic}
\setmainfont{Linux Libertine O}[Ligatures=TeX]
\newfontfamily\arabicfont{Amiri}[Script=Arabic]
\newfontfamily\ArNote{Amiri}[Script=Arabic,Scale=1.05]
\hypersetup{pdftitle={al-Battani, al-Zij al-Sabi: an English translation of the canons},pdfauthor={al-Battani}}
\setlength{\parindent}{0pt}\setlength{\parskip}{0.55em}
\tracinglostchars=2
\newcommand{\ArPage}[1]{\marginpar{\footnotesize\color{gray}Ar.\,p.\,#1}}
\newcommand{\ChapterHead}[2]{\par\bigskip\penalty-300\pdfbookmark[1]{#1}{ch-#2}\hypertarget{EN-C#2}{}%
 {\large\bfseries #1}\par\medskip}
\begin{document}

\thispagestyle{empty}
\begin{center}\Large AL-BATTĀNĪ\par\vspace{4mm}\LARGE The Sabian Zīj\par\vspace{2mm}\large \emph{al-Zīj al-Ṣābiʾ}\par\vspace{5mm}\normalsize the canons, translated from the Arabic\par\vspace{2mm}in Nallino's edition, Part III (Rome, 1899)\end{center}\vspace{8mm}
\noindent This is an English translation of the canons of al-Battānī's \emph{Zīj}, made from the Arabic text as Nallino printed it. Every sentence of the Arabic is translated, in the order of the text; the numbers are given in figures. Words added for the sense stand in square brackets. The notes identify quotations, record where the translation follows a correction of Nallino's, and explain terms. The pages of the Arabic text are given in the margin. Nallino's Latin translation (Part I, Milan, 1903) was used as a check on the meaning throughout; its wording is not followed.
\par\smallskip\noindent The translation is in progress. This version holds the list of chapters and chapters 1, 2. It is AI-integrated work (Claude), 2026; no person has proofread it. It is dedicated to the public domain (CC0 1.0).
\par\clearpage
\begin{center}\ArPage{1}In the name of God, the Merciful, the Compassionate. God bless the Prophet Muḥammad, His noble Messenger, and his family and his companions, and grant them peace.\end{center}
\par This is the whole of what Muḥammad ibn Jābir ibn Sinān al-Ḥarrānī, known as al-Battānī — may God pardon him — set down on the computation of the stars and on the places of their courses as they have been tested. The book contains fifty-seven chapters in all, and this is an account of them.\footnote{The heading gives the author's name as «Muḥammad ibn Jābir ibn Sinān»; the title page that Nallino set above it has «ibn Sinān ibn Jābir». «Tested» (\emph{al-mumtaḥan}) is the word al-Battānī uses for what has been verified by observation.}\par
\par\noindent\hangindent=2em\hangafter=1 \makebox[2em][l]{1.}The opening of the book.\par
\par\noindent\hangindent=2em\hangafter=1 \makebox[2em][l]{2.}On dividing the circle of the sphere, and on multiplying degrees by one another, taking their roots, and dividing them by one another.\par
\par\noindent\hangindent=2em\hangafter=1 \makebox[2em][l]{3.}On finding the lengths of the chords of the parts of the circle, on entering in tables the halves of the chords of double the arcs, and on the procedures with them that follow from this.\footnote{The halves of the chords of double the arcs are the sines; chapter 3 introduces them under this name.}\par
\par\noindent\hangindent=2em\hangafter=1 \makebox[2em][l]{4.}On the amount of the inclination of the ecliptic to the equator; on the division of this inclination, its directions, and its stages as it rises and descends. It is the declination of the sun from the right sphere.\footnote{The print has «\RL{\ArNote ونجزية}»; Nallino's corrigenda read «\RL{\ArNote وتجزية}», «the division», which is translated.}\par
\par\noindent\hangindent=2em\hangafter=1 \makebox[2em][l]{5.}On finding the amounts of the equator that rise together with given arcs of the ecliptic [in places] under the equator, which is called the line of equality. By these same amounts, too, the signs pass over the meridian circle in every place on the earth; this is why they are called the ascensions of the signs in the right sphere.\footnote{Nallino suspects that the words «in the countries that are» have fallen out before «under the equator»; the bracket supplies their sense. «The line of equality» (\emph{khaṭṭ al-istiwāʾ}) is the terrestrial equator.}\par
\par\noindent\hangindent=2em\hangafter=1 \ArPage{2}\makebox[2em][l]{6.}On the properties of each of the lines parallel to the equator that lie away from it towards the north, and an account of the places of the inhabited earth whose longitude and latitude are known from the book \emph{The Form of the Earth}.\par
\par\noindent\hangindent=2em\hangafter=1 \makebox[2em][l]{7.}On finding the amplitudes of the winter and summer risings and settings on the horizon circles of the countries. These are the arcs of the horizon between the equator and the places of the ecliptic, and they are called the azimuths of the risings and settings on the horizon circle.\par
\par\noindent\hangindent=2em\hangafter=1 \makebox[2em][l]{8.}On finding the altitude of the north pole from the excess of the longest day, when this is given.\par
\par\noindent\hangindent=2em\hangafter=1 \makebox[2em][l]{9.}On finding the excess of the longest day from the given altitude of the pole.\par
\par\noindent\hangindent=2em\hangafter=1 \makebox[2em][l]{10.}On finding the altitude and the shadow, each from the other, when the shadow is horizontal, and on finding them when the shadow is vertical.\par
\par\noindent\hangindent=2em\hangafter=1 \makebox[2em][l]{11.}On finding the azimuth of the altitude and of the shadow on the horizon circle, in every country, at every time of the day, for all degrees of the ecliptic. This is the part of the horizon circle, counted from the limit of rising and setting, that is cut off by the arc passing through the zenith and the sun.\par
\par\noindent\hangindent=2em\hangafter=1 \makebox[2em][l]{12.}On finding the meridian line in every country, which is the direction of the south, and the direction of the equinoctial rising that becomes known with it; and on finding it in various ways.\par
\par\noindent\hangindent=2em\hangafter=1 \makebox[2em][l]{13.}On finding the amount of the equator that rises together with given degrees of the ecliptic from the horizon in every place on the earth, which is called the ascensions of the signs in each country; and what follows from this: finding the ascensions of any portion you wish in these ascensions and in the ascensions of the right sphere, finding the degrees of the ecliptic from these ascensions, the amount of the arc of the day and of the night, their equal hours, the times of the seasonal hours of the day and of the night, and converting the one kind of hours into the other.\footnote{Nallino notes that Plato of Tivoli's copy seems to have had «degree» (\emph{daraja}) for «portion» (\emph{wajh}), which may be the better reading.}\par
\par\noindent\hangindent=2em\hangafter=1 \makebox[2em][l]{14.}On finding the latitudes of the countries, which is the altitude of the north pole above their horizon, by observation.\par
\par\noindent\hangindent=2em\hangafter=1 \makebox[2em][l]{15.}On finding the altitude of the sun at noon on any day.\par
\par\noindent\hangindent=2em\hangafter=1 \makebox[2em][l]{16.}On finding how many hours of the day have passed and what is rising, by measuring the sun and knowing the altitude and the vertical shadow.\par
\par\noindent\hangindent=2em\hangafter=1 \makebox[2em][l]{17.}On finding the altitude from the hours of the day that have passed.\par
\par\noindent\hangindent=2em\hangafter=1 \ArPage{3}\makebox[2em][l]{18.}On finding the distances of the fixed or the wandering stars from the equator when they lie away from the belt of the signs in latitude, and the degrees of the ecliptic that culminate with them, from their places in longitude and latitude with respect to the ecliptic.\par
\par\noindent\hangindent=2em\hangafter=1 \makebox[2em][l]{19.}On finding the half of the arc of the day of a star, which is half of its stay above the earth, and also below it, and the times of its hours above the earth and below it.\footnote{«The half» is supplied by Nallino; the codex and Plato have «the arc of the day».}\par
\par\noindent\hangindent=2em\hangafter=1 \makebox[2em][l]{20.}On finding the degree of the ecliptic with which a star rises and the degree with which it sets.\par
\par\noindent\hangindent=2em\hangafter=1 \makebox[2em][l]{21.}On finding how many hours of the night have passed, by measuring some of the stars.\par
\par\noindent\hangindent=2em\hangafter=1 \makebox[2em][l]{22.}On finding the altitude of some of the stars from the hours of the night that have passed.\par
\par\noindent\hangindent=2em\hangafter=1 \makebox[2em][l]{23.}On finding the azimuth of a star from its altitude above the horizon.\par
\par\noindent\hangindent=2em\hangafter=1 \makebox[2em][l]{24.}On finding the distance of a star from the equator, and the parts of the ecliptic that culminate with it, from the azimuth of the place on the horizon circle where it rises or sets. By this the declination of a degree of the ecliptic from the equator is also known.\footnote{Al-Battānī distinguishes the distance of a star from the equator, when the star lies off the ecliptic, from the declination (\emph{mayl}) of the degrees of the ecliptic and of the sun; so the last sentence of the title does not repeat the first.}\par
\par\noindent\hangindent=2em\hangafter=1 \makebox[2em][l]{25.}On finding the degree of the ecliptic in which a star is, and the latitude of the star, from its distance from the equator and the degree that culminates with it, when these are known.\par
\par\noindent\hangindent=2em\hangafter=1 \makebox[2em][l]{26.}On finding the distances between the stars according to the arrangement of their places on the sphere in longitude and latitude.\par
\par\noindent\hangindent=2em\hangafter=1 \makebox[2em][l]{27.}On finding the length of the times of the solar year as found by observation, and from it the mean motion of the sun in days, months and years.\par
\par\noindent\hangindent=2em\hangafter=1 \makebox[2em][l]{28.}On finding the irregularity of the motion of the sun, and what appears with it of the degrees of the signs from the places of its farthest distance.\footnote{The Arabic is as ambiguous as Plato's Latin. Nallino takes the sense to be the longitude of the sun reckoned from its apogee, as it comes out from the irregular motion of the sun.}\par
\par\noindent\hangindent=2em\hangafter=1 \makebox[2em][l]{29.}On finding the amounts of the difference between days with their nights, when the daytime of one day together with its night is compared with the daytime of another day together with its night, and how they are converted and transferred from one to the other.\par
\par\noindent\hangindent=2em\hangafter=1 \ArPage{4}\makebox[2em][l]{30.}On the description of the spheres of the moon and its motions; the irregularity that appears in them at the times of the conjunctions and oppositions with the sun, and the second irregularity that is compounded with it on account of the moon's distances from the sun; the causes of the two kinds of eclipse; the distance of the two luminaries from the earth; and the waxing and waning of the moon's light with its distance from the sun.\par
\par\noindent\hangindent=2em\hangafter=1 \makebox[2em][l]{31.}On the description of the spheres of the wandering stars and their states.\par
\par\noindent\hangindent=2em\hangafter=1 \makebox[2em][l]{32.}On finding the eras of the Arabs, the Rūm, the Persians and the Copts, and finding each of them from another.\par
\par\noindent\hangindent=2em\hangafter=1 \makebox[2em][l]{33.}On finding the place of the ecliptic in which the sun is seen, by the era of the Rūm or of the Arabs, whichever of the two you wish.\par
\par\noindent\hangindent=2em\hangafter=1 \makebox[2em][l]{34.}On finding the hours of computation in every country. These are the mean equal hours after noon at the city of al-Raqqa, by which the motions are found at any time, so that the mean position of a star is known at that time of the day or the night; and on converting these hours into the hours of the countries.\par
\par\noindent\hangindent=2em\hangafter=1 \makebox[2em][l]{35.}On setting up the ascendant and the twelve houses from the hours of the day and the night, and finding the hours from the ascendant.\par
\par\noindent\hangindent=2em\hangafter=1 \makebox[2em][l]{36.}On finding the true place of the moon on the ecliptic on any day and at any time.\par
\par\noindent\hangindent=2em\hangafter=1 \makebox[2em][l]{37.}On finding the place of the northern and the southern node, which are the head and the tail, through which the moon passes in latitude.\par
\par\noindent\hangindent=2em\hangafter=1 \makebox[2em][l]{38.}On finding the latitude of the moon, which is its distance from the belt of the signs towards the south and the north.\par
\par\noindent\hangindent=2em\hangafter=1 \makebox[2em][l]{39.}On finding the parallax that the moon undergoes in longitude and latitude, its amounts in the quarters of the horizon, and the cause from which it arises; in various ways.\par
\par\noindent\hangindent=2em\hangafter=1 \makebox[2em][l]{40.}On finding the distance of the moon from the earth from its parallax in the circle of altitude that lies between the zenith and the horizon and passes through the place of the moon on the ecliptic.\par
\par\noindent\hangindent=2em\hangafter=1 \makebox[2em][l]{41.}On the sighting of the crescent at the beginning and the end of the months; the azimuth of the place in the sky where it is seen; its altitude above the horizon at that time; and its form, according to the light that is in it and to whether its two horns are level or tilted.\par
\par\noindent\hangindent=2em\hangafter=1 \makebox[2em][l]{42.}On finding the mean and the true conjunctions and oppositions of the sun and the moon, by the era of the Rūm or of the Copts, whichever you wish.\par
\par\noindent\hangindent=2em\hangafter=1 \makebox[2em][l]{43.}On finding the lunar eclipses, their amounts and their times in the countries; the side of the moon's disc on which the eclipse begins and the side on which the clearing takes place; and its figure, and working it out by computation and by the tables.\par
\par\noindent\hangindent=2em\hangafter=1 \ArPage{5}\makebox[2em][l]{44.}On finding the eclipse of the sun, its different amounts in every country and its times there; the side of the sun's disc on which the eclipse begins and the side on which it clears; and its figure, and working it out by computation and by the tables.\par
\par\noindent\hangindent=2em\hangafter=1 \makebox[2em][l]{45.}On finding the places of the five wandering stars on the ecliptic at any time.\par
\par\noindent\hangindent=2em\hangafter=1 \makebox[2em][l]{46.}On finding the stations of the five wandering stars and their retrogradation.\par
\par\noindent\hangindent=2em\hangafter=1 \makebox[2em][l]{47.}On finding the latitudes of the five wandering stars and their directions.\par
\par\noindent\hangindent=2em\hangafter=1 \makebox[2em][l]{48.}On finding the appearance and the disappearance of the five wandering stars.\par
\par\noindent\hangindent=2em\hangafter=1 \makebox[2em][l]{49.}On finding the nine configurations that the fixed stars and some of the wandering stars have with respect to the sun.\par
\par\noindent\hangindent=2em\hangafter=1 \makebox[2em][l]{50.}An account of the distances of the stars from the earth, their diameters, the sizes of their bodies and the extent of their spheres.\par
\par\noindent\hangindent=2em\hangafter=1 \makebox[2em][l]{51.}On finding the motion of the other stars by observation, and on entering in the table the places of those that are needed, in longitude and latitude.\footnote{Where this list has «the other stars», the heading of the chapter itself and Plato's Latin have «the fixed stars».}\par
\par\noindent\hangindent=2em\hangafter=1 \makebox[2em][l]{52.}On what the masters of talismans have said, that the sphere has a motion of displacement going forwards and backwards, and the error that appears in this.\footnote{This is the theory later known as trepidation, an oscillation of the equinoxes. Nallino translates «the masters of talismans» (\emph{aṣḥāb al-ṭilasmāt}) as «the astrologers».}\par
\par\noindent\hangindent=2em\hangafter=1 \makebox[2em][l]{53.}On finding the times of the revolutions of the years, which take place when the sun returns to the place in which it was at the beginning.\par
\par\noindent\hangindent=2em\hangafter=1 \makebox[2em][l]{54.}On determining the amounts of the applications according to the latitudes of the stars, when they cast their rays on the ecliptic.\footnote{The applications (\emph{ittiṣālāt}) are the astrological aspects between stars.}\par
\par\noindent\hangindent=2em\hangafter=1 \makebox[2em][l]{55.}On finding the ascensions of the signs between the quadrants of the sphere.\par
\par\noindent\hangindent=2em\hangafter=1 \makebox[2em][l]{56.}On making the vertical and the horizontal sundial for finding the seasonal hours of the day in every country; on setting it up correctly and on the direction of the south; and how the direction of the qibla, which is the direction of Mecca the protected, is found on the sundial.\footnote{The sundial is \emph{al-rukhāma}, literally «the marble slab».}\par
\par\noindent\hangindent=2em\hangafter=1 \makebox[2em][l]{57.}The conclusion of the book, and a description of the making of the instrument in the shape of the sphere, which is called «the egg», and of the making of the two instruments for observation.\footnote{Nallino suspects «making» (\emph{ṣanʿa}) where the codex has \emph{ḍaʿa}, «placing»; Plato has «making». The translation follows him.}\par
\par This is the account of the contents of the book in detail, which comprises fifty-seven sections. Praise be to God for His help, and God bless Muḥammad.\par
\ChapterHead{Chapter 1. The opening of the book}{01}\ArPage{6}
\par He said: The first thing with which every matter is begun and every discourse is opened is the praise of God, exalted be His mention, thanks to Him for His benefits, and the blessing on the Seal of His messengers and prophets — upon them be peace, and the mercy of God and His blessings.\footnote{The Seal of the messengers and prophets is Muḥammad (cf. Qurʾān 33:40).}\par
\par Praise be to God, who created the creatures by His power, ordered all things by His will and perfected them by His wisdom; who «encompasses all things in knowledge» and «has counted all things in number»; from whom «not the weight of an atom in the heavens or on the earth escapes, nor anything smaller than that or greater, but it is in a clear book».\footnote{Qurʾān 65:12, 72:28 and 34:3. The «clear book» is the book that is with God.}\par
\par I bear witness that there is no god but God alone, who has no partner, and I bear witness that Muḥammad is His servant and His messenger, whom «He sent with the guidance and the religion of truth, to make it prevail over all religion, though the idolaters dislike it». Through him He guided the believers and «cut off the last of the unbelievers», and He made him a proof for all the worlds. God bless him and his good family, his chosen companions and those who follow his practice, until the Day of Judgement.\footnote{Qurʾān 9:33 and 61:9 (Nallino's note has XLI, 9, a misprint for LXI, 9); for «cut off the last of the unbelievers», cf. 7:72 and 8:7.}\par
\par To proceed. After the knowledge of the laws and the practice of religion, of which no one may be ignorant, the science of the art of the stars is among the sciences that are highest in rank and most exalted in standing, finest in adornment, that cling most to the heart and shine most in the soul, and that do most to sharpen thought and reflection, to quicken understanding and to train the mind. For there is a great share of good in it and great benefit: the knowledge of the length of the years and months, of the appointed times, of the seasons, of the lengthening and shortening of the day and the night, of the places of the two luminaries and their eclipses, of the courses of the stars in their direct and retrograde motion, of the changes in their configurations, of the order of their spheres and of all their other relations; and beyond this, what is attained through it by one who examines it closely and reflects on it constantly: the proof of the unity of God, and the knowledge of the essence of the Creator's greatness, of the breadth of His wisdom, the majesty of His power and the subtlety of His handiwork.\footnote{The configurations of the stars are the nine configurations with the sun that chapter 49 explains (Nallino).}\par
\par \ArPage{7}He has said — mighty is the Speaker: «In the creation of the heavens and the earth, and in the alternation of night and day, there are signs for those who have understanding.» He, blessed and exalted, has said: «Blessed is He who has placed signs in the sky.» He, mighty and majestic, has said: «It is He who has made the night and the day to follow one another.» He, glory be to Him, has said: «It is He who has made the sun a radiance and the moon a light, and has measured out its stations, so that you may know the number of the years and the reckoning.» And He, exalted be His mention, has said: «The sun and the moon run by a reckoning.» There is much more of this kind in the book of God, mighty and majestic, which it would take long to describe and much space to mention and cite.\footnote{Qurʾān 3:190, 25:61, 25:62, 10:5 and 55:5. The signs (\emph{burūj}) of 25:61 are the word that al-Battānī uses for the signs of the zodiac throughout the book. In the third quotation the text prints «the night and the moon a light, in succession», a slip that mixes this verse with the next quotation; the Qurʾān (25:62) has «the night and the day», which Nallino translates and the translation follows.}\par
\par When I had studied this science at length and reflected on it constantly, I came to know how the books composed on the motions of the stars disagree, and into what errors some of their authors fell in the procedures they set down in them as foundations and in what they built on these. I also came to know what has accumulated in the motions of the stars over long periods of time, when their observations are compared with the ancient observations; that the inclination of the ecliptic to the equator has been found to be smaller; and what has changed with this change in the various kinds of computation, in the lengths of the years, in the times of the seasons, and in the conjunctions of the two luminaries that are inferred from the times and hours of the eclipses. So I set about correcting all this and making it exact according to the method of Ptolemy in the book known as the Almagest, after close study, long reflection and consideration, following in his footsteps and keeping to what he laid down. For he had investigated these matters from every side and had shown their causes and the reasons that occur in them by a geometrical and numerical proof whose soundness cannot be rejected and whose truth cannot be doubted. After that he ordered that they be tested and examined after him, and said that his observations might be improved upon in the course of time, as he himself had improved upon Hipparchus and others of his peers; for the art is a lofty one, and, being heavenly and vast, it can be grasped only approximately.\footnote{«Has been found to be smaller» renders \emph{mā wujida … min al-taqārub}, literally «the drawing together that was found»: al-Battānī's observations gave a smaller obliquity of the ecliptic than the Greek value (chapter 4). Nallino stresses that the words do not show knowledge of a continuous decrease of the obliquity. For Ptolemy's remark see the Almagest, III 1 (in Heiberg's numbering; Nallino cites Halma's III, 2).}\par
\par I have composed a book on this, in which I have made clear what was obscure, opened what was closed and explained what was difficult in the principles of this science and out of the way in its branches. With it I have made the path of guidance easy for those who follow it and work by it in the art of the stars. In it I have corrected the motions of the stars and their places on the belt of the ecliptic, as I found them by observation and by the computation of the two kinds of eclipse, together with all the other procedures that are needed, and I have added other things that are needed. I have arranged that the motions of the stars are found from the tables for noon of the day for which the computation is made, at the city of al-Raqqa, where the observation and the testing that establish all this took place — if God, exalted is He, wills. Success comes from God.\par
\ChapterHead{Chapter 2. On the division of the circle of the sphere, and on multiplication, roots and division}{02}\ArPage{8}
\par He said: The ancients divided the circle of the sphere into 360 parts, and they gave more than one reason for this. One is that the number of these parts is close to the number of days of the year, which is completed when the sun passes a fixed point of the sphere until it returns to it. Another is that it is a number that has a half, a third, a quarter and other fractions that do not come out whole for many numbers. They placed the sun at four points of the sphere, which bring about two equinoxes and two solstices and divide the year into four distinct parts, spring, summer, autumn and winter; and they named each of these points after the season that begins when the sun passes it.\par
\par Since everything that has extension has a middle and two ends, each of these seasons divides into three parts, and so the circle of the sphere had to have twelve parts. They found the spring point the best of these points and the most fitting to begin with. For from it, after the equinox, the day begins to lengthen, and the sun is rising into the northern half of its sphere, so that the heat grows; and the nature of this season is moist and inclines to heat, which suits the beginning of growth and the coming into being of things. So they made the reckoning of the sphere begin from it.\footnote{Nallino notes the same account of the spring point in Abū Maʿshar (d. 886), \emph{Introductorium} II 5, and, almost word for word, in Abraham bar Ḥiyya's \emph{Sphaera mundi} (1165), ch. 2.}\par
\par Then they found that the constellations next to these twelve parts, which are called signs, were twelve constellations, and they named each sign after the constellation next to it, although over a long time these constellations move away from the places of the signs named after them. So the first part became Aries, followed by Taurus, Gemini, Cancer, Leo, Virgo, Libra, Scorpio, Sagittarius, Capricorn, Aquarius and Pisces.\footnote{The Arabic names are al-Ḥamal (the Ram), al-Thawr (the Bull), al-Jawzāʾ (Gemini), al-Saraṭān (the Crab), al-Asad (the Lion), al-Sunbula (the Ear of Grain), al-Mīzān (the Balance), al-ʿAqrab (the Scorpion), al-Qaws (the Bow), al-Jady (the Kid), al-Dalw (the Bucket) and al-Ḥūt (the Fish). The translation uses the Latin names by which the same signs are known in English.}\par
\par Each of these signs therefore has 30 parts as its share of the 360 parts of the circle of the sphere. These parts are also called degrees. Each degree is divided into 60 parts, called minutes; each minute is likewise divided into 60 parts, called seconds; each second is divided into 60 thirds; and so on, by the same rule of division, to tenths and the separate orders that come after them.\par
\par \ArPage{9}As for the meaning of multiplication: it is repeating one of two numbers as many times as there are units in the other — I mean the multiplication of units by units. Multiplying fractions by units is repeating the fractions as many times as there are units, or taking of the units the part that the fractions are of one. Multiplying fractions by fractions is taking of one of the two fractions, whichever you wish, the part that the other fraction is of one.\par
\par Thus, when degrees are multiplied by degrees, the product is degrees; when they are multiplied by minutes, it is minutes; when they are multiplied by seconds, the product is seconds; and so with what is multiplied by thirds, fourths and the following orders: the product is of the order of the smaller one, by which it was multiplied. As for the orders below degrees, minutes and the others: when each of them is multiplied by itself, the product falls below it by as much as it falls below degrees. For example, minutes multiplied by minutes give seconds, and multiplied by seconds they give thirds, and what is multiplied by thirds and fourths falls by the same rule. Seconds multiplied by seconds give fourths, and multiplied by thirds they give fifths; and whatever lies further from these orders goes the same way, by the same rule.\par
\par When any number of one of these orders, obtained by multiplication or by addition, is divided by sixty, the number at which every order ends and to which all the fractions are related, the result belongs to the next higher order. When, in a number made up of two or more of these orders, more has to be subtracted from one of them than it contains, one unit is broken off from the next higher order for it, counted as 60 parts and added to it; then as much as needed is subtracted, and what remains is counted together with what remains of the higher order.\par
\par As for degrees: whatever accumulates of them when the parts of the motions are added is reckoned in revolutions. If the sum is more than one revolution, or than several revolutions, a revolution being 360 degrees, the revolutions are dropped and what remains is counted. When more has to be subtracted from the degrees than they contain, a revolution is added to them; then as much as needed is subtracted from the sum, and what remains is counted.\par
\par \ArPage{10}If you wish to multiply one of the orders of degrees or of fractions by another of them, and to know from this table of which order the product will be, take on one of the two lines AB the cell in which the order you wish to multiply by any other order is written, and go straight on from that cell until you are level with the other order that you want: the cell level with it, among the orders of fractions, gives the order of the product. For example, suppose you wish to multiply fourths by thirds. Take on the line AB that runs across the width of the page either of the two orders — let it be the thirds first — and go on from it until you are level with the fourths on the line AB that runs along the length of the page; in the cell level with it you find sevenths, and this is the order of the product. Likewise, if you took the fourths on a line AB and went from them to the level of the thirds on the other line AB, you would find sevenths there. Do the same with any orders you wish, if God wills.\par
\begin{landscape}\begin{center}{\bfseries The orders of products}\footnote{The print reads from right to left; here the columns are set from left to right. The first row is the line AB across the width of the page, the first column the line AB along its length. The cell of the fourth row and second column is printed «\RL{\ArNote ورابع}», a misprint for «\RL{\ArNote روابع}», «fourths».}\par\medskip{\footnotesize\setlength{\tabcolsep}{3pt}\begin{tabular}{|c|c|c|c|c|c|c|c|c|c|c|}\hline degrees & minutes & seconds & thirds & fourths & fifths & sixths & sevenths & eighths & ninths & tenths \\
\hline minutes & seconds & thirds & fourths & fifths & sixths & sevenths & eighths & ninths & tenths & elevenths \\
\hline seconds & thirds & fourths & fifths & sixths & sevenths & eighths & ninths & tenths & elevenths & twelfths \\
\hline thirds & fourths & fifths & sixths & sevenths & eighths & ninths & tenths & elevenths & twelfths & thirteenths \\
\hline fourths & fifths & sixths & sevenths & eighths & ninths & tenths & elevenths & twelfths & thirteenths & fourteenths \\
\hline fifths & sixths & sevenths & eighths & ninths & tenths & elevenths & twelfths & thirteenths & fourteenths & fifteenths \\
\hline sixths & sevenths & eighths & ninths & tenths & elevenths & twelfths & thirteenths & fourteenths & fifteenths & sixteenths \\
\hline sevenths & eighths & ninths & tenths & elevenths & twelfths & thirteenths & fourteenths & fifteenths & sixteenths & seventeenths \\
\hline eighths & ninths & tenths & elevenths & twelfths & thirteenths & fourteenths & fifteenths & sixteenths & seventeenths & eighteenths \\
\hline ninths & tenths & elevenths & twelfths & thirteenths & fourteenths & fifteenths & sixteenths & seventeenths & eighteenths & nineteenths \\
\hline tenths & elevenths & twelfths & thirteenths & fourteenths & fifteenths & sixteenths & seventeenths & eighteenths & nineteenths & twentieths \\\hline\end{tabular}}\end{center}\end{landscape}
\par \ArPage{11}As for the meaning of the root: the root of any absolute number, whatever number it is, is that which, multiplied by itself, gives the given number. But the roots of these orders do not keep this rule, because of the difference that we have described in what results from multiplying these orders by one another. The rule holds only for the order of degrees: the root of degrees is also degrees, since degrees multiplied by degrees give degrees. As for the fractions below degrees, of the remaining orders: those of even rank, such as seconds, fourths, sixths and the like, have their root in the higher order whose rank is half of theirs — seconds, whose root is minutes, and fourths, whose root is seconds. Those of odd rank, such as minutes, thirds and the like, have no definite root unless they are converted into the order below them, so that they come to an even rank; then this rule applies to them: minutes are converted into seconds, and thirds into fourths.\footnote{An absolute number (\emph{ʿadad muṭlaq}) is a number taken by itself, not counted in one of the sexagesimal orders.}\par
\par Division is finding how many times the smaller number goes into the larger, when the larger is measured by the smaller, and finding what part of the larger the smaller is, when the smaller is the number divided. Here we proceed in the reverse of what we did in multiplications and roots, by the same rule: if we divide degrees by degrees, the quotient is degrees. As for the other orders below degrees: when a lower order is divided by a higher one, whatever its rank and whether or not it is next to it, the quotient is of the order which, multiplied by the order divided by, gives back the order of the number divided. Thus seconds divided by minutes give minutes, and sixths divided by fourths give seconds. But when a higher order is divided by a lower one, the way to do it is to convert the higher order into the lower and then divide by it, and the quotient is degrees. So with minutes divided by sixths: when they are converted into sixths and then divided by those sixths, the result of the division is degrees, as we have described.\par
\par \ArPage{12}If you wish to know, with the table given here, what results from dividing the lower orders of fractions by the orders above them, look up in the line AB or in the line AD, whichever you wish, the order that you wish to divide by an order of higher rank, whether or not it is next to it, and go on level with it until you are level with the higher order on the other line. The order of fractions at which you arrive is the order that you obtain by dividing the number divided, and it is the order which, multiplied by the higher order that you divided by, gives back the order of the number divided. Likewise, if you wish to divide a higher order by a lower one, convert the higher into the lower, look up in one of the two lines the order into which it has been converted, and go on level with it until you are level with the order that you wish to divide by: you obtain degrees. And whenever you divide an order by the same order, the result is degrees, if God, exalted is He, wills. Success comes from God.\footnote{«With the table given here»: the Arabic says «with this table mentioned before»; the table is printed at the head of the next page.}\par
\begin{landscape}\begin{center}{\bfseries The orders of quotients}\footnote{The columns are set from left to right, as in the first table. The first row is the line AB, the first column the line AD. The cells of the second row and fourth column and of the fourth row and second column are printed «\RL{\ArNote ثولنٍ}» and «\RL{\ArNote توانٍ}», misprints for «\RL{\ArNote ثوانٍ}», «seconds».}\par\medskip{\footnotesize\setlength{\tabcolsep}{3pt}\begin{tabular}{|c|c|c|c|c|c|c|c|c|c|c|}\hline degrees & minutes & seconds & thirds & fourths & fifths & sixths & sevenths & eighths & ninths & tenths \\
\hline minutes & degrees & minutes & seconds & thirds & fourths & fifths & sixths & sevenths & eighths & ninths \\
\hline seconds & minutes & degrees & minutes & seconds & thirds & fourths & fifths & sixths & sevenths & eighths \\
\hline thirds & seconds & minutes & degrees & minutes & seconds & thirds & fourths & fifths & sixths & sevenths \\
\hline fourths & thirds & seconds & minutes & degrees & minutes & seconds & thirds & fourths & fifths & sixths \\
\hline fifths & fourths & thirds & seconds & minutes & degrees & minutes & seconds & thirds & fourths & fifths \\
\hline sixths & fifths & fourths & thirds & seconds & minutes & degrees & minutes & seconds & thirds & fourths \\
\hline sevenths & sixths & fifths & fourths & thirds & seconds & minutes & degrees & minutes & seconds & thirds \\
\hline eighths & sevenths & sixths & fifths & fourths & thirds & seconds & minutes & degrees & minutes & seconds \\
\hline ninths & eighths & sevenths & sixths & fifths & fourths & thirds & seconds & minutes & degrees & minutes \\
\hline tenths & ninths & eighths & sevenths & sixths & fifths & fourths & thirds & seconds & minutes & degrees \\\hline\end{tabular}}\end{center}\end{landscape}
\end{document}
